\documentclass[aspectratio=1610,svgnames]{beamer}

\usepackage{lmodern}
\usepackage[T1]{fontenc}
\usepackage[ngerman]{babel}
\usepackage{selinput}
\SelectInputMappings{%
   adieresis={ä},
   germandbls={ß}
   }

\usetheme{PaloAlto}  %% Themenwahl

\setbeamercovered{transparent}
%\setbeamertemplate{footline}[frame number]
\usecolortheme{spruce}		% grün
\usecolortheme[named=MSUgreen]{structure}
\newcommand{\hshblogo}{\includegraphics[width=1.3cm]{space-logo}}
\newcommand{\divider}[1]{\begin{frame} %
\begin{alertblock}{} %
\centering\usebeamerfont{section title}#1 %
\end{alertblock} %
\end{frame}}
 
\title{vim --- Eine Einführung}
\author{Thomas Helmke}
\date{02.12.2014}
\logo{\includegraphics[width=1.1cm]{space-logo}}
 
\begin{document}
\maketitle
\frame{\tableofcontents}

\section{Einleitung}
\divider{\insertsection}
\begin{frame}[<+->] %%Eine Folie
	\frametitle{Warum vim?} %%Folientitel
	\begin{itemize}
		\item Warum nicht?
		\item Effizientes Arbeiten (nach der Lernkurve)
		\item Ohne Maus bedienbar
	\end{itemize}
\end{frame}
\begin{frame}[<+->] %%Eine Folie
	\frametitle{Unterschied zu ``normalen'' Editoren} %%Folientitel
	\begin{itemize}
		\item Verschiedene Modi für unterschiedliche Zwecke
		\item Eigener Kommandosatz
		\item Für reine Tastaturbedienung konzipiert
	\end{itemize}
\end{frame}

\section[Die Modi]{Die verschiedenen Modi}
\divider{\insertsection}
\begin{frame}[<+->]
    \frametitle{Normal-Modus}
    \begin{itemize}
        \item Der Grundmodus
        \item Zum Navigieren durch die Datei
        \item Mit <ESC> kommt man zurück
    \end{itemize}
\end{frame} 
\begin{frame}[<+->]
    \frametitle{Visual-Modus}
    \begin{itemize}
        \item Zum Text markieren
        \item Zeichenweise, Zeilenweise, Blockweise
        \item Hier funktioniert die Maus
    \end{itemize}
\end{frame} 
\begin{frame}[<+->]
    \frametitle{Insert-Modus}
    \begin{itemize}
        \item Erreichbar durch i, a, o, \ldots
        \item Hier kann man endlich Text eingeben
        \item Nutzung wird geringer je mehr man von vim lernt
    \end{itemize}
\end{frame} 

\end{document}
